\section{Counting global words and anchored local patterns}\label{sec:counting}

All counts in this section are counts of literal interval labels.
The central observation is that each scalar coordinate can be counted
using an arrangement in two real parameters, even though the number
$D$ of coordinates grows with $k$.
We use the standard hyperplane-arrangement region count; see
Stanley~\cite[Lemma~2.1 and Proposition~2.4]{stanley2007}.
The elementary proof below also counts equality signs, which are needed
for the half-open mesh intervals.

\begin{lemma}[Affine comparisons, including equality]\label{count:arrangement}
Let $s$ affine hyperplanes be given in $\R^e$, and record, for each
hyperplane, the sign in $\{-,0,+\}$ of a defining affine function.
The number of realized joint sign vectors is at most
\begin{equation}\label{count:arrangement-bound}
 \sum_{d=0}^e\binom{s}{d}\sum_{j=0}^{e-d}\binom{s}{j}
 \le(e+1)^2(s+1)^e.
\end{equation}
The same bound applies when the parameter space is restricted to an
arbitrary subset of $\R^e$.
\end{lemma}

\begin{proof}
First count strict sign vectors. The complement of an arrangement
has at most $\sum_{j=0}^e\binom{s}{j}$ regions. Indeed, when a new
hyperplane is added, every region it splits meets that hyperplane in
a different region of the induced arrangement, whose dimension is
at most $e-1$. Induction on $s$ and $e$, followed by Pascal's identity,
gives the stated bound. Repeated hyperplanes create no additional
regions. A strict sign set is convex, and therefore connected when
nonempty, so each strict sign vector corresponds to one region.

Now consider a realized sign vector with zero entries, and let $A$
be the affine intersection of its zero hyperplanes. If $A$ has
codimension $d$, then $d$ of those equations with linearly independent
normals already define $A$: consistency at a point of $A$ makes every
remaining zero equation a consequence of the selected equations.
There are at most $\binom{s}{d}$ ways to select them. On $A$, every
hyperplane containing $A$ has forced zero sign; the remaining
hyperplanes impose strict signs in dimension $e-d$. The strict-region
bound on $A$, and a sum over $d$, prove the first inequality in
\Cref{count:arrangement-bound}. For each summand,
$\binom{s}{d}\binom{s}{j}\le(s+1)^{d+j}\le(s+1)^e$;
there are at most $(e+1)^2$ summands. Restricting the parameter space
can only remove sign vectors.
\end{proof}

\begin{lemma}[Global label words]\label{count:global}
Let $T_{\mathrm{glob}}$ be the number of words
$\bigl(L(n_j)\bigr)_{0\le j<k}$ arising from the progressions in
\Cref{setup:cyclic}. Then
\begin{equation}\label{count:global-bound}
 \log T_{\mathrm{glob}}
 =O\bigl(D(\log k+\log\lambda)\bigr)
 =O(k^{3/10}\log k)=o(M).
\end{equation}
\end{lemma}

\begin{proof}
Consider any circle partition with $F$ endpoints, including its
representative cut, and freely choose a scalar start and step
$a,b\in[0,1)$. For $0\le j<k$, the unreduced value $a+jb$ belongs
to $[0,k)$. Represent the partition endpoints by numbers
$\tau\in[0,1)$. The interval containing $(a+jb)\bmod1$ is determined
by its comparisons with all $\tau+z$ for $-1\le z\le k$.
The comparisons for all $j$ use at most $3k^2F$ lines in the
two parameters $a,b$. Equality signs give precisely the assigned
half-open labels, so \Cref{count:arrangement} bounds the number of
scalar words by $9(3k^2F+1)^2$.

Apply this bound separately in all $2D$ scalar coordinates, allowing
their starts and steps to vary independently. This enlarges the
collection of possible words, and is therefore an upper bound for
the actual coordinate homomorphisms. We can take
\[
 F=|U|+|V|=k^2\lambda+O(k^3\log k).
\]
Since $\lambda\ge1$, we have
$\log F=O(\log k+\log\lambda)$, and thus
\[
 T_{\mathrm{glob}}\le\bigl[9(3k^2F+1)^2\bigr]^{2D}.
\]
This proves the first bound in \Cref{count:global-bound} with an
absolute constant. The remaining bounds follow from
\Cref{setup:orders} and $M\ge k^{1/2}$.
\end{proof}

For the local patterns fix an integer $h$ with $h_0<h\le2M$ and
$t\in\{0,\ldots,h-1\}^D$. In coordinate $i$ let
\begin{equation}\label{count:coordinate-periods}
 s_i=\frac{h}{\gcd(h,\lambda t_i)}.
\end{equation}
Call that second-system coordinate \emph{stationary} if $s_i=1$, and
\emph{rotating} otherwise. By \Cref{setup:dilation}, a rotating
coordinate has $s_i>h_0$. These descriptions depend only on $h,t$.
Stationarity concerns the rational step and permits nonzero drift;
even a period $h>h_0$ may have only stationary second coordinates.

We now define an enlarged family of full return-period blocks whose
labels can be balanced before the outer coloring is chosen.
\Cref{sec:dichotomy} will verify membership for every actual full
block with $h>h_0$ containing a position whose full label occurs more
than $k/M$ times in the whole progression.

\begin{definition}[Eligible local patterns]\label{count:eligible}
Choose $x,y\in\T^D$ and $u,v\in\R^D$, and form
\begin{equation}\label{count:local-paths}
 X_z=x+\frac z h(t+u),\qquad
 Y_z=y+\frac z h(\lambda t+v)\qquad(0\le z<h)
 \quad\text{in }\T^D.
\end{equation}
A position $z$ is \emph{regular} if
$\rho(Y_{z,i})\ge\eta$ in every rotating coordinate; if there are
no rotating coordinates, every position is regular. Record
$L(X_z,Y_z)$ at a regular position and a dummy symbol $*$ otherwise.
The resulting word in $(\cL\cup\{*\})^h$ is \emph{eligible for
$h,t$} if it admits a realization satisfying all of the following:
\begin{enumerate}[label=\textup{(\roman*)}]
 \item At least $h/2$ positions are regular, and the full labels at
       those positions are pairwise distinct.
 \item $|u_i|\le H_x$ for every $i$, and $|v_i|\le H$ in every
       rotating coordinate.
 \item In every stationary coordinate $i$, one has
       $|v_i|\le |V(Y_{z,i})|$ at every regular position $z$.
\end{enumerate}
The recorded word is the pattern; different realizations of one word
do not define additional patterns.
\end{definition}

The stationary drift bound is relative to interval width because those
intervals may be arbitrarily fine. This is the condition that makes
the following count independent of the anchor's width.

\begin{lemma}[Local patterns through a specified label]\label{count:local}
There is an absolute constant $C$ such that, for every
$h_0<h\le2M$ and every full label $\beta\in\cL$, the number of pairs
consisting of $t\in\{0,\ldots,h-1\}^D$ and an eligible pattern for
$h,t$ containing $\beta$ is at most
\begin{equation}\label{count:local-bound}
 (CDh)^{6D}.
\end{equation}
The constant is independent of $q$, $\lambda$, $\beta$, and $D$.
\end{lemma}

\begin{proof}
Fix $t$ and a position $z_0$ at which the recorded pattern equals
$\beta$. That position is regular. In each scalar coordinate,
re-anchor the path at $z_0$: let $a$ be the real representative of
that coordinate at $z_0$, in the interval specified by $\beta$,
and let $b$ be its drift variable, either $u_i$ or $v_i$.
At a target position $z$, an unreduced representative is
\begin{equation}\label{count:anchored-scalar}
 a+\tau_z+\frac{z-z_0}{h}b,
 \qquad
 \tau_z=\frac{(z-z_0)t_i}{h}
 \quad\hbox{or}\quad
 \tau_z=\frac{(z-z_0)\lambda t_i}{h},
\end{equation}
as appropriate. The translation $\tau_z$ is fixed once $t,z_0,z,i$
are fixed. We retain only the restrictions on $a,b$ implied at the
anchor, since discarding other restrictions can only increase the
number of outputs. The following three cases bound this number.

\emph{First-system coordinates.} Here $a$ ranges over an interval of length
$H_x$ and $|b|\le H_x$. Since $|z-z_0|/h<1$, the target range
is contained in a fixed real interval of length at most $3H_x$.
The real lifts of $U$-endpoints form a lattice of spacing $H_x$,
so at most four endpoints meet this range.

\emph{Stationary second-system coordinates.} Write $I=V(Y_{z_0,i})$.
The retained drift bound is $|b|\le |I|$, by the restriction at
the regular anchor. Also $\lambda t_i/h$ is integral, so the
translation $\tau_z$ is integral. Modulo one every target belongs
to $N(I)$ from \Cref{mesh:neighbors}, and at most $49$ mesh endpoints
are relevant. The unreduced target range lies in an interval of
length at most $3|I|<1$, so each relevant circle endpoint has at
most one lift in that range.

\emph{Rotating second-system coordinates.} The anchor interval has
width at most $2H$ and the retained drift bound is $|b|\le H$.
The target range has length at most $4H<1$. Apply
\Cref{mesh:truncated} to its image on the circle: at most $16004D$
breakpoints determine the truncated output $V_\eta$. Again each
circle breakpoint has at most one lift in the real target range.

In each case, the enclosing target interval and the finite list of
breakpoint lifts are fixed before $a,b$ are varied. In particular,
the magnitude of the fixed translation $\tau_z$ introduces no
factor in the count. Each comparison with a breakpoint lift is an
affine line in $(a,b)$. Taken over $0\le z<h$, there are at most
$16004Dh$ such lines per scalar coordinate. On the allowed parameter
rectangle, their signs determine the scalar output at every target.
They determine values at breakpoints as well as on open intervals;
this accounts for half-open endpoints and for exact equalities
$\rho=\eta$.

By \Cref{count:arrangement} with $e=2$, the number of possible scalar
output sequences is at most
\[
 9(16004Dh+1)^2\le C_0(Dh)^2
\]
for an absolute $C_0$. Take the product of these bounds for all $2D$
coordinates, truncating only rotating second-system coordinates.
Those scalar outputs determine the complete recorded pattern: a
position is dummy exactly when at least one rotating output is $*$;
at all other positions the outputs specify the full label. No
additional choice of a regular-position mask is needed.

There are $h^D$ choices of $t$ and at most $h$ choices of the anchor
$z_0$. Therefore the required number is at most
\[
 h^{D+1}\bigl[C_0(Dh)^2\bigr]^{2D}
 =C_0^{2D}D^{4D}h^{5D+1}
 \le(CDh)^{6D},
\]
where the last inequality holds for an absolute $C$, since $D\ge1$.
All eligible patterns containing $\beta$ have at least one such
anchor and are included in this count.
\end{proof}
